% 同一环境分阶段接入两种策略；右侧与静态图共用分布、采样、训练和评价的读法。
\begin{tikzpicture}[
  box/.style={draw=NNLine,line width=1.1pt,rounded corners=2pt,
    text width=3.7cm,minimum height=1.1cm,align=center,
    font=\figurefont\small,inner sep=5pt},
  data/.style={box,draw=NNInput,fill=FigureInputFill},
  train/.style={box,draw=NNHidden,fill=FigureRedFill},
  result/.style={box,draw=NNOutput,fill=FigureYellowFill},
  flow/.style={-{Stealth[length=2mm]},draw=NNWire,line width=1.2pt},
  analysis/.style={flow,dashed},
  annotation/.style={font=\figurefont\footnotesize,align=center},
  every node/.append style={text=NNInk}]
  \draw[draw=NNLine,line width=.85pt,rounded corners=3pt,fill=NNPanel]
    (-2.15,-6.25) rectangle (2.15,.9);
  \node[annotation,anchor=south west] at (-2.15,1.0) {策略与同一环境的交互};
  \node[result] (beta) at (0,0) {采集：行为策略$\beta$\\$a_t\sim\beta_t(\cdot\mid I_t)$};
  \node[data,minimum height=1.45cm] (env) at (0,-2.7) {同一环境$\mathcal E$\\$\mu_0,T$；观测核$K_o$};
  \node[result] (pi) at (0,-5.4) {执行：学得策略$\pi$\\$a_t\sim\pi_t(\cdot\mid I_t)$};
  \draw[flow] ([xshift=-.65cm]beta.south) --
    node[left,annotation] {行动$a_t$} ([xshift=-.65cm]env.north);
  \draw[flow] ([xshift=.65cm]env.north) --
    node[right,annotation] {$o_t,I_t$} ([xshift=.65cm]beta.south);
  \draw[flow] ([xshift=-.65cm]pi.north) --
    node[left,annotation] {行动$a_t$} ([xshift=-.65cm]env.south);
  \draw[flow] ([xshift=.65cm]env.south) --
    node[right,annotation] {$o_t,I_t$} ([xshift=.65cm]pi.north);

  \node[data] (train-law) at (5.15,0) {轨迹律与记录分布\\
    $P_{\mathcal E}^{\beta}$\\$P_{\mathrm{rec}}^{\beta}=b_{\#}P_{\mathcal E}^{\beta}$};
  \node[data] (sample) at (10.3,0) {可见训练轨迹样本\\$D\sim(P_{\mathrm{rec}}^{\beta})^n$};
  \draw[flow] (beta) -- node[above,annotation] {诱导} (train-law);
  \draw[flow] (train-law) -- node[above,annotation] {独立\\采样} (sample);

  \node[train] (objective) at (5.15,-2.7) {训练目标与可行集\\$F_D$及训练约束};
  \node[result] (learned) at (10.3,-2.7) {学得策略\\$\pi=A(D,U)$};
  \draw[flow] (sample.south) -- (10.3,-1.25)
    -- node[above,annotation] {构造$B$} (5.15,-1.25) -- (objective.north);
  \draw[flow] (objective) -- node[above,annotation] {优化} (learned);
  \draw[flow] (learned.south) -- (10.3,-3.55)
    -- node[below,annotation] {固定学得策略后执行} (1.8,-3.55)
    -- ([xshift=1.8cm]pi.north);

  \node[data] (test-law) at (5.15,-5.4) {轨迹律与记录分布\\
    $P_{\mathcal E}^{\pi}$\\$P_{\mathrm{rec}}^{\pi}=b_{\#}P_{\mathcal E}^{\pi}$};
  \node[data] (test-sample) at (10.3,-5.4) {可见测试轨迹样本\\$D_{\mathrm{te}}\sim(P_{\mathrm{rec}}^{\pi})^m$};
  \draw[flow] (pi) -- node[above,annotation] {诱导} (test-law);
  \draw[flow] (test-law) -- node[above,annotation] {独立\\采样} (test-sample);

  \node[result] (risk) at (5.15,-7.4) {总体轨迹风险$R_{\mathcal M}(\pi)$\\
    $\mathbb E_{P_{\mathcal E}^{\pi}}[\ell(\pi,\tau)]$};
  \node[result] (estimate) at (10.3,-7.4) {风险估计\\$\widehat R_{D_{\mathrm{te}}}(\pi)$};
  \draw[analysis] (test-law) -- node[left,annotation] {期望定义} (risk);
  \draw[flow] (test-sample) -- (estimate);
  \draw[flow] (learned.east) -- (13,-2.7) -- (13,-7.4) -- (estimate.east);
  \draw[analysis] (estimate) -- node[above,annotation] {统计\\分析} (risk);
\end{tikzpicture}
